\section{Discussion}
\label{sec:nt-discussion}

\paragraph{Branding as fashion.}
\citet{simmel1957fashion} describes fashion with two forces: imitation, which spreads a style, and differentiation, which motivates it. A coined name is a differentiation device: it marks a contribution as a distinct artefact. Once 40\% of submissions carry one, the device no longer differentiates. The diffusion we document, a logistic rise strongest in fast-growing topics and among newcomer-containing teams, is the imitation half of that cycle and has the S-shape of classic diffusion \citep{bass1969new}. Abrahamson's management-fashion theory adds that technical fashions must look rational and progressive \citep{abrahamson1991managerial,abrahamson1996management}. A name does that work: it presents a paper as a reusable component, a ``transportable package'' in \citeauthor{fujimura1988bandwagon}'s (\citeyear{fujimura1988bandwagon}) sense. Our decomposition supports a mixed reading. Depending on design, a third to two thirds of the rise tracks what papers are about; in fast-growing topics branding runs ahead of what bag-of-words content predicts, even before LLMs; and in the one additive split more than half stays unexplained by measured content, composition or detectable tool use. Internal dynamics of taste \citep{lieberson2000matter} or informational cascades \citep{bikhchandani1992theory} are consistent with that residual, but we cannot distinguish them from a neutral-drift baseline \citep{bentley2004random} or from unmeasured content. Fast-rising tastes are predicted to be abandoned faster \citep{berger2009adoption}; the \NAME{} share is still rising in 2026, so whether a downswing follows \citep{abrahamson1999lifecycles} is open and testable.

\paragraph{Communal versus generational change.}
For co-authored text the sociolinguistic distinction between speakers who change and cohorts that replace each other is harder to draw than it looks, because every co-author shares the same title. Our evidence favours change carried by new keys and new projects: the gradient in team newness, the absence of net \NAME{} additions within retitled resubmission pairs, and the post-2023 first-position newcomer premium. It also shows an asymmetry. Incumbent keys track topic shifts fully but style shifts only partly ($G_{\text{style}}$ 0.17 \ci{0.13}{0.20}; 0.40 at the first list position, 0.01 at the last). This fits the finding that half of 2023 LLM-paper first authors were new to LLM research \citep{movva2024topics}, and suggests that norms of scientific presentation turn over partly through cohort replacement even when research agendas turn over within incumbent keys.

\paragraph{Implications for meta-science.}
\begin{itemize}
\item \emph{Include all submissions.} Withdrawn and desk-rejected papers are part of the population in which fashion arises. A published-only study reproduces 9 of 11 mainstream statistics within $\pm 10\%$ (\NAME{} 1.021), but its point estimate of the 2024 \NAME{} slope change is half the all-submission value (0.50 \ci{-0.15}{1.16}; pool equality not rejected, $p=0.51$), so it may miss part of the acceleration. It under-represents outsider styles (ALL-CAPS 0.36 \ci{0.26}{0.47}), which the audit cannot separate from camera-ready recasing.
\item \emph{Treat post-decision snapshots as measurement problems.} Comparing pre-decision arXiv versions with tipping points and a multiplicity-controlled gate is cheap and reusable. Here it turned an apparent decision-stage premium into an unidentified one.
\item \emph{Use year-level inference} when the treatment varies by cohort, and \emph{use a generational null} for any key-level statistic computed on co-authored text.
\end{itemize}

\paragraph{Implications for LLM-era writing studies.}
Excess-word methods \citep{kobak2025delving,liang2025quantifying} measure topic and style together. Meaning-fixed slots help separate them for short texts, conditional on a slot being filled: in 2024 the new verbs were added to titles that opened with a verb rather than swapped for synonyms, and substitution appears only in 2025--26. The link of this wording to tool use runs through abstract fingerprints of unknown recall, and fingerprint-based tool shares are model-conditional; in our case the fingerprint behind the tool share failed its own positive control. Fixed-vocabulary bounds also decay as model vocabulary changes. Finally, nothing here shows that review penalises LLM text: the one new LLM-marker depletion is small (3--4\%), sits before the decision, and concerns abstracts, not titles.
